\documentclass[12pt]{article}
\usepackage[utf8]{inputenc}
\usepackage[T1]{fontenc}
\usepackage{hyperref}
\title{Temb\'e Language Resources for K-12 Curriculum Development}
\author{Compiled by Claude Sonnet 5 (claude-sonnet-5) for Andr\'as Kornai}
\date{}
\begin{document}
\maketitle

\section{Introduction}
Temb\'e (ISO 639-3 tqb) is a Tupian language of Par\'a and Maranh\~ao in Brazil. It is a variety of the Tenetehara complex, which also includes Guajajara. The two are mutually intelligible but are treated as separate languages by their communities. One source says fewer than a quarter of the roughly 820 Temb\'e people spoke the language in 2000. Another says that among the Temb\'e near the Guam\'a River it is no longer spoken.

\section{Primary Educational Resources}
No school materials for Temb\'e turned up in the searches.

\section{Linguistic and Reference Materials}
\begin{itemize}
\item Max Boudin, \emph{Dicion\'ario de Tupi moderno: dialeto temb\'e-t\'en\^et\'ehar do alto rio Gurupi}, S\~ao Paulo, 1978. It is cited in bibliographies, and I did not find an online copy.
\item F\'abio Bonfim Duarte's theses on Temb\'e clause structure (UnB, 1997) and constituent order (UFMG, 2003). The 2003 thesis may be the file at \url{https://repositorio.ufmg.br//bitstreams/34f8e8b6-3909-47df-b3c7-6b78ea82a6d3/download}, though the search result gave no title for that file.
\item Tabita Fernandes da Silva, 2010 UnB dissertation on the history of the Tenetehara complex: \url{https://repositorio.unb.br/bitstream/10482/8947/3/2010_TabitaFernandesdaSilva.pdf}
\item Wikipedia, Tenetehara language: \url{https://en.wikipedia.org/wiki/Tenetehara_language}
\item Povos Ind\'igenas no Brasil, Temb\'e: \url{https://www.indios.org.br/en/Povo:Temb%C3%A9}
\end{itemize}

\section{Religious and Community Resources}
None were found.

\section{Digital and Audio Resources}
None were found.

\section{Curriculum Development Context}
Scholarly descriptions exist, but I found no pedagogical materials and no recordings. Any curriculum work should start with the Tenetehara/Guajajara literature and with Temb\'e community speakers. I did not verify whether Guajajara school materials could be adapted.
\end{document}